% concatenated from 20260515.tex
\documentclass[12pt,reqno]{amsart}

\usepackage{amsmath,amssymb,amsthm,mathrsfs}
\usepackage{graphicx,cite,times}
\usepackage{cases}
\usepackage{xcolor}
\usepackage{appendix}
\usepackage{enumerate}
\usepackage{soul}
\usepackage[colorlinks, linkcolor=blue, citecolor=blue]{hyperref}

\usepackage{xpatch}
\usepackage{hyperref}
\usepackage[final]{showkeys}

\soulregister{\cite}7





\setlength{\topmargin}{-1.5cm}
\setlength{\oddsidemargin}{0.0cm}
\setlength{\evensidemargin}{0.0cm}
\setlength{\textwidth}{16.7cm}
\setlength{\textheight}{23cm}
\headheight 20pt
\headsep    26pt
\footskip 0.4in

\newtheorem{theo}{Theorem}[section]
\newtheorem{coro}[theo]{Corollary}
\newtheorem{lemm}[theo]{Lemma}
\newtheorem{prop}[theo]{Proposition}
\newtheorem{rema}[theo]{Remark}
\newtheorem{defi}[theo]{Definition}
\newtheorem{nota}[theo]{Notation}

\numberwithin{equation}{section}




\title[Inverse Problem of Biharmonic wave  Equations]{ Lipschitz Stability for an Inverse Problem of Biharmonic Wave Equations with Damping
}

\author{  Minghui Bi  }
\address{School of Mathematics and Statistics, Center for Mathematics and
	Interdisciplinary Sciences, Northeast Normal University, Changchun, Jilin 130024, P.R.China. }
\email{bimh801@nenu.edu.cn}



\author{Yixian Gao}
\address{School of Mathematics and Statistics, Center for Mathematics and
	Interdisciplinary Sciences, Northeast Normal University, Changchun, Jilin 130024, P.R.China. }
\email{gaoyx643@nenu.edu.cn}



\thanks{ The research of Y. Gao was supported by NSFC grant
	12371187 and STDP Project of Jilin Province
	20240101006JJ}




\subjclass[2010]{35R30, 35L35, 35B35}


\keywords{ Biharmonic Wave Equations, Parameter Identification,  Inverse Problem, Stability}




\begin{document}
	
	\begin{abstract}
		We consider the inverse problem of the simultaneous identification of a variable density coefficient and the initial state for a damped biharmonic wave equation. Given boundary Cauchy data of the Laplacian of the state, we prove that the parameters are determined with Lipschitz stability. The analysis begins by showing that the underlying evolution constitutes a contraction semigroup, establishing the well-posedness of the direct problem. By employing multiplier methods, we derive a crucial observability inequality that serves as the cornerstone for our stability analysis. We obtain explicit stability estimates that characterize the interplay between the higher-order operator structure and the regularity of the coefficients. Notably, we quantify the influence of the damping mechanism on the reconstruction, showing that the stability constants scale with the factor $(1+\gamma)^{1/2}$. These results provide a rigorous theoretical framework for the study of high-order hyperbolic inverse problems.
		
		
		
		
		
	\end{abstract}
	
	\maketitle
	
	\section{Introduction\label{sect1}}
	
	\subsection{Statement of the main results}
	
	The biharmonic wave equation serves as a canonical model for the transverse vibrations of thin elastic plates. In this paper, we are concerned with the inverse problem of the simultaneous identification of internal material properties—specifically a spatially heterogeneous density—and the initial state of the system from boundary measurements.
	
	We consider the damped biharmonic wave equation with variable density $\rho(x)$ on a bounded domain $\Omega \subset \mathbb{R}^n$ with a smooth boundary $\partial\Omega$:
	\begin{align}\label{a}
		\begin{cases}
			\rho(x) \partial_t^2 u(t, x) + \Delta^2 u(t, x) + \gamma \partial_t u(t, x) = 0, & (t, x) \in \mathbb{R}_+ \times \Omega, \\[6pt]
			u(0, x) = f(x), \quad \partial_t u(0, x) = g(x), & x \in \Omega, \\[6pt]
			u(t, x) = 0, \quad \partial_n u(t, x) = 0, & (t, x) \in \mathbb{R}_+ \times \partial \Omega.
		\end{cases}
	\end{align}
	Here, $\gamma \ge 0$ is the damping coefficient and $\Delta^2$ denotes the biharmonic operator. The initial displacement and velocity are given by $f$ and $g$, respectively. Our framework accommodates density functions $\rho(x)$ that may be piecewise constant; a prototypical example (cf. \cite{Faouzi2025}) is the configuration
	\[
	\rho(x) = \rho_0 + (\rho_1 - \rho_0) \mathbf{1}_{\omega}(x),
	\]
	where $\omega \Subset \Omega$ represents a star-shaped inclusion and $\rho_0, \rho_1 > 0$ are constants.
	
	The inverse problem consists of recovering the unknown density $\rho$ and the initial displacement $f$ from the boundary Cauchy data of the Laplacian of the state:
	\[
	\Delta u|_{\partial \Omega} \quad \text{and} \quad \partial_n(\Delta u)|_{\partial \Omega},
	\]
	recorded over a finite time interval $[0, T]$. Our primary objective is to establish Lipschitz stability for this identification. To distinguish the density variation from the acceleration term—which effectively acts as a source in the linearized problem—we impose a standard non-degeneracy condition on the acceleration of the state. Physically, this ensures the system is not in a steady state or a regime of vanishing acceleration.
	
	We define the class of admissible parameters as follows. The density $\rho$ is a measurable function satisfying 
	\[ 0 < \rho_{\min} \le \rho(x) \le \rho_{\max} < \infty \quad \text{a.e. } x \in \Omega. \]
	The initial displacement $f$ is assumed to belong to $H^{4}(\Omega) \cap H^{2}_{0}(\Omega)$ and satisfies the compatibility condition $\Delta^{2}f|_{\partial\Omega}=0$, which ensures $H^4$-regularity of the solution. The initial velocity $g$ belongs to $H^{2}_{0}(\Omega)$. Any triple $(\rho, f, g)$ satisfying these regularity and boundary conditions is termed \emph{admissible}.
	
	Our first main result establishes the Lipschitz stability of the density coefficient.
	
	\begin{theo}\label{th1}
		Let $u_{1}$ and $u_{2}$ be the solutions of \eqref{a} corresponding to two admissible triples $(\rho_{1},f_{1},g_{1})$ and $(\rho_{2},f_{2},g_{2})$, and let $u := u_{1} - u_{2}$. 
		Suppose the acceleration of the second solution satisfies the uniform non-degeneracy condition
		\[
		\inf_{t\in[0,T]} \|\partial_t^2 u_2(t,\cdot)\|_{L^2(\Omega)} \ge \alpha > 0
		\]
		and the regularity bound $\|\partial_t^2 u_2\|_{L^{\infty}(0,T;L^{2}(\Omega))} \le M$. Then there exists a constant $C > 0$, depending only on $(M,T,\alpha,\rho_{\min},\rho_{\max},\Omega)$, such that
		\begin{align}\label{eq:stability-estimate}
			\|\rho_{1}-\rho_{2}\|_{L^{\infty}(\Omega)}
			\le C\,(1+\gamma)^{1/2} \left( \int_{0}^{T}\int_{\partial\Omega} \left( \left| \partial_n \Delta u \right|^{2} + |\Delta u|^{2} \right) \mathrm{d} S \mathrm{d}t \right)^{1/2}.
		\end{align}
	\end{theo}
	
	Under the additional assumption that the initial velocities coincide, the boundary data also determines the initial displacement with the same stability scaling.
	
	\begin{theo}\label{th2}
		Let the hypotheses of Theorem \ref{th1} hold and assume $g_{1}=g_{2}$. If the observation time $T$ satisfies the geometric condition
		\[
		T > \frac{2\,\mathrm{diam}(\Omega)}{\sqrt{\rho_{\min}}},
		\]
		then there exists a constant $C > 0$ such that
		\[
		\|f_{1}-f_{2}\|_{H^{4}(\Omega)}
		\le C\,(1+\gamma)^{1/2} \left( \int_{0}^{T}\int_{\partial\Omega} \left( \left| \partial_n \Delta u \right|^{2} + |\Delta u|^{2} \right) \mathrm{d}S \mathrm{d}t \right)^{1/2}.
		\]
	\end{theo}
	
	\noindent
	These estimates demonstrate that the biharmonic structure provides a robust mechanism for parameter recovery. The stability constants exhibit an explicit square-root dependence on the damping parameter $\gamma$, quantifying the influence of energy dissipation on the conditioning of the inverse problem.
	
	\subsection{Background and Motivation}
	
	The mathematical analysis of hyperbolic equations traces its origins to foundational studies of the early twentieth century. These early works focused on the construction of fundamental solutions, the solvability of singular boundary value problems, and the influence of dissipative mechanisms on well-posedness \cite{MR1577376, MR1577103, MR3908, MR28496, MR21214}. Together, these investigations established the requisite operator-theoretic framework for the subsequent development of higher-order evolution models. 
	
	The study of biharmonic wave equations, in particular, arose from classical elasticity theory and boundary value problems for thin plates. Early contributions by Mackevi\v{c} and Atakhodzhaev \cite{Atahodzaev1974} utilized elliptic coordinate systems and regularization techniques to establish existence and uniqueness for biharmonic systems in complex geometries. Subsequent developments, such as the work of Arzhanykh \cite{Arzhanykh1978} on quasi-analytic solutions and Ramm \cite{Ramm1988} on the half-plane Cauchy problem, further expanded the analytical scope of biharmonic models.
	
	In recent years, attention has increasingly shifted toward the inverse problem of coefficient identification, progressing from foundational uniqueness results to sharp stability estimates and nonlinear generalizations. Contemporary progress frequently relies on complex geometrical optics (CGO) solutions and time-frequency analysis. For instance, stability estimates for density inversion in damped fourth-order equations were derived in \cite{MR4598840, Liu2025FourthOrderSchrodinger}, achieving the simultaneous reconstruction of the density and internal sources. In the nonlinear regime, high-frequency transformation methods have been applied to Westervelt-type models \cite{Acosta2025Westervelt}. Furthermore, recent advances have relaxed structural assumptions on the coefficients \cite{MR4959941}, extending the theory to one-dimensional evolution equations \cite{Kian2025WaveEquation} and unbounded domains \cite{MR4848562}.

Despite these advances, the simultaneous identification of multiple parameters in 
higher-order hyperbolic systems remains a significant challenge. At the foundational 
theoretical level, Li et al.\ have successfully addressed the qualitative uniqueness 
problem for biharmonic inverse scattering. They established the unique continuation 
principle for fourth-order biharmonic wave equations with rough potentials and 
systematically developed a unified framework for uniqueness analysis in both 
deterministic and random media \cite{li2026sampling, li2026stability, wang2025adaptive, 
harris2025direct, li2025nonradiating, dong2024uniqueness, dong2024novel}. Furthermore, 
the microlocal analysis and high-frequency asymptotic techniques introduced in these 
works have become standard tools for theoretical and computational research in this field.

While qualitative uniqueness has been largely resolved, quantitative stability and 
multi-parameter inversion pose persistent open problems. The higher-order nature of the 
biharmonic operator complicates the derivation of the Carleman and observability 
estimates necessary for Lipschitz stability; consequently, classical methods developed 
for lower-order or simplified models \cite{Benrabah2019, Danh2019} cannot be directly 
extended to fourth-order systems. Moreover, existing literature often relies on 
restrictive assumptions: the high-frequency stability analysis in \cite{Zhao2025} 
neglects the coupling effects between spatially varying density and damping, while the 
random source inversion in \cite{Li2022a} fails to quantify how regularity loss impacts 
reconstruction stability. Finally, most existing results are restricted to single-parameter 
inversion or simplified geometries---for instance, \cite{Bouslah2024} focuses exclusively 
on boundary reconstruction in simply connected domains, and \cite{Bhattacharyya2025} does 
not address simultaneous parameter recovery. Therefore, the joint reconstruction of 
density and initial states, alongside the explicit quantification of damping effects, 
remains an unresolved challenge.

	
	
	To address these gaps, the present work introduces a unified framework for the simultaneous reconstruction of the medium density and the initial displacement. We first establish the well-posedness of the forward problem by proving that the underlying evolution operator generates a contraction semigroup. A central contribution of this paper is the derivation of a novel energy observability inequality via multiplier techniques, which serves as the cornerstone for our inverse analysis. Furthermore, we provide the first explicit Lipschitz stability estimate for this coupled problem, characterizing the precise scaling of the stability constant with respect to the damping parameter $\gamma$. These results confirm that the biharmonic structure inherently enhances the conditioning of the inverse problem, thereby providing a rigorous theoretical basis for applications in non-destructive evaluation and dynamic inversion.
	
	The remainder of the paper is organized as follows. Section \ref{sect2} addresses the well-posedness of the forward problem within the natural energy space $\mathcal{H}$. In Section \ref{sec3}, we derive the fundamental observability inequality. Section \ref{sec4} is devoted to the proofs of Theorems \ref{th1} and \ref{th2}, which establish the Lipschitz stability estimates for the density and initial displacement. In Section \ref{sec5}, we present numerical simulations in both one- and two-dimensional configurations to empirically verify the theoretical stability estimates, specifically illustrating the dependence of the stability constant on the damping parameter. Finally, Section \ref{sec6} provides concluding remarks.
	
	
	
	
	
	
	
	
	
	
	
	
	\section{Properties of the operator \( \mathcal{A} \) \label{sect2}}
	
	To establish the well-posedness of the forward problem \eqref{a}, we recast the system as an abstract Cauchy problem in an appropriate Hilbert space. We define the energy space
	\[
	\mathcal{H} := \left( H^4(\Omega) \cap H_0^2(\Omega) \right) \times H_0^2(\Omega),
	\]
	equipped with the inner product
	\[
	\left\langle \begin{pmatrix} u_1 \\ v_1 \end{pmatrix}, \begin{pmatrix} u_2 \\ v_2 \end{pmatrix} \right\rangle_{\mathcal{H}} = \int_\Omega \Delta^2 u_1  \overline{\Delta^2 u_2} \, \mathrm{d}x + \int_\Omega \rho \, v_1 \overline{v_2} \, \mathrm{d}x,
	\]
	which induces a norm equivalent to the standard $H^4 \times H^2$ norm. The second term represents a weighted $L^2$-inner product with respect to the density $\rho \in L^\infty(\Omega)$.
	
	We define the system operator $\mathcal{A} : D(\mathcal{A}) \subset \mathcal{H} \to \mathcal{H}$ by
	\[
	\mathcal{A} \begin{pmatrix} u \\ v \end{pmatrix} = \begin{pmatrix} v \\ -\rho^{-1} \Delta^2 u - \rho^{-1} \gamma v \end{pmatrix},
	\]
	with the domain
	\[
	D(\mathcal{A}) = \Bigl\{ (u, v) \in \mathcal{H} \ \Big|\ \Delta^2 u + \gamma v \in \rho H_0^2(\Omega) \Bigr\}.
	\]
	The well-posedness of \eqref{a} follows from the maximal dissipativity of $\mathcal{A}$. We first establish that $\mathcal{A}$ generates a contraction semigroup on $\mathcal{H}$, ensuring the existence of a unique solution for admissible initial data.
	
	\begin{lemm}[Dissipativity] \label{lemm1}
		The operator $\mathcal{A}$ is dissipative. That is, for every $(u, v) \in D(\mathcal{A})$, 
		\[
		\operatorname{Re} \left\langle \mathcal{A} \begin{pmatrix} u \\ v \end{pmatrix}, \begin{pmatrix} u \\ v \end{pmatrix} \right\rangle_{\mathcal{H}} \leq 0.
		\]
	\end{lemm}
	
	\begin{proof}
		Let $(u, v) \in D(\mathcal{A})$. By the definition of the inner product in $\mathcal{H}$, we have
		\[
		\begin{aligned}
			\left\langle \mathcal{A} \begin{pmatrix} u \\ v \end{pmatrix}, \begin{pmatrix} u \\ v \end{pmatrix} \right\rangle_{\mathcal{H}} 
			&= \int_\Omega \Delta^2 v \, \overline{\Delta^2 u} \, \mathrm{d}x + \int_\Omega \rho \left( -\rho^{-1} \Delta^2 u - \rho^{-1} \gamma v \right) \overline{v} \, \mathrm{d}x \\[4pt]
			&= \int_\Omega \Delta^2 v \, \overline{\Delta^2 u} \, \mathrm{d}x - \int_\Omega \Delta^2 u \, \overline{v} \, \mathrm{d}x - \gamma \int_\Omega |v|^2 \, \mathrm{d}x.
		\end{aligned}
		\]
		Consider the first two terms of the expression. Since $\Delta^2$ is self-adjoint on $H^4(\Omega) \cap H_0^2(\Omega)$, and noting that $v \in H_0^2(\Omega)$ and $u \in H^4(\Omega) \cap H_0^2(\Omega)$, an application of Green's identity (or the definition of the biharmonic operator in the sense of distributions) yields:
		\[
		\int_\Omega \Delta^2 v \, \overline{\Delta^2 u} \, \mathrm{d}x - \int_\Omega \Delta^2 u \, \overline{v} \, \mathrm{d}x = \int_\Omega \Delta^2 v \, \overline{\Delta^2 u} \, \mathrm{d}x - \int_\Omega \Delta u \, \overline{\Delta v} \, \mathrm{d}x.
		\]
		For $u \in H^4$ and $v \in H^2$, the term $I = \int_\Omega \Delta^2 v \overline{\Delta^2 u} \, \mathrm{d}x - \int_\Omega \Delta^2 u \overline{v} \, \mathrm{d}x$ satisfies $\operatorname{Re}(I) = 0$ due to the anti-symmetry of the undamped part of the evolution in this specific energy topology. Specifically,
		\[
		\operatorname{Re} \left\langle \mathcal{A} \begin{pmatrix} u \\ v \end{pmatrix}, \begin{pmatrix} u \\ v \end{pmatrix} \right\rangle_{\mathcal{H}} = -\gamma \int_\Omega |v|^2 \, \mathrm{d}x.
		\]
		Given that $\gamma \geq 0$, the result follows.
	\end{proof}
	
	
	By the Lumer--Phillips theorem, the generation of a $C_0$-semigroup of contractions by a dissipative operator $\mathcal{A}$ on a Hilbert space $\mathcal{H}$ is equivalent to the surjectivity of the map $(\lambda I - \mathcal{A})$ for some $\lambda > 0$. We now establish this property for the system operator.
	
	\begin{lemm}[Maximality] \label{lemm2} 
		The operator $\mathcal{A}$ is maximally dissipative.
	\end{lemm}	
	
	\begin{proof}
		In view of Lemma \ref{lemm1}, it suffices to show that $\operatorname{ran}(\lambda I - \mathcal{A}) = \mathcal{H}$ for any $\lambda > 0$. Given an arbitrary $\mathbf{y} = (y_1, y_2) \in \mathcal{H}$, we seek $\mathbf{x} = (u, v) \in D(\mathcal{A})$ such that $(\lambda I - \mathcal{A})\mathbf{x} = \mathbf{y}$. This resolvent equation is equivalent to the coupled system
		\begin{align}
			\lambda u - v &= y_1, \label{eq:res1} \\
			\rho \lambda v + \Delta^2 u + \gamma v &= \rho y_2. \label{eq:res2}
		\end{align}
		Substituting $v = \lambda u - y_1$ from \eqref{eq:res1} into \eqref{eq:res2}, we obtain the following fourth-order elliptic equation for $u$:
		\begin{equation}\label{eq:elliptic_final}
			\Delta^2 u + (\lambda^2 \rho + \lambda \gamma) u = \rho y_2 + (\lambda \rho + \gamma) y_1 \quad \text{in } \Omega,
		\end{equation}
		supplemented with the clamped boundary conditions $u = \partial_n u = 0$ on $\partial\Omega$. 
		
		Since $y_1 \in H^4(\Omega) \cap H_0^2(\Omega)$ and $y_2 \in H_0^2(\Omega)$, the right-hand side of \eqref{eq:elliptic_final} belongs to $L^2(\Omega)$. The coefficient function $a(x) := \lambda^2 \rho(x) + \lambda \gamma$ is strictly positive and resides in $L^\infty(\Omega)$. By the Lax--Milgram theorem and standard elliptic regularity for the biharmonic operator with Dirichlet conditions, there exists a unique solution $u \in H^4(\Omega) \cap H_0^2(\Omega)$ to \eqref{eq:elliptic_final}.
		
		Defining $v := \lambda u - y_1$, it follows immediately that $v \in H_0^2(\Omega)$ because $y_1 \in H_0^2(\Omega)$ and $u$ inherits this regularity. Finally, from \eqref{eq:res2}, we have
		\[
		-\rho^{-1} \Delta^2 u - \rho^{-1} \gamma v = \lambda v - y_2.
		\]
		Since $v \in H_0^2(\Omega)$ and $y_2 \in H_0^2(\Omega)$, the right-hand side is in $L^2(\Omega)$, which confirms that $\mathbf{x} = (u, v) \in D(\mathcal{A})$. This establishes the surjectivity of $\lambda I - \mathcal{A}$ and completes the proof.
	\end{proof}
	
	
	Although maximal dissipativity implies closedness in this setting, we provide a direct verification for completeness.
	
	\begin{lemm}[Closedness] \label{lemm3}
		The operator \( \mathcal{A} \) is closed in \( \mathcal{H} \).
	\end{lemm}
	
	\begin{proof}
		Let \( \{\mathbf{x}_n\}_{n\in\mathbb{N}} = \{(u_n, v_n)\}_{n\in\mathbb{N}} \subset D(\mathcal{A}) \) be a sequence such that 
		\begin{equation*}\label{eq:conv_sequences}
			\mathbf{x}_n \xrightarrow{\mathcal{H}} \mathbf{x} = (u, v) \quad \text{and} \quad \mathcal{A}\mathbf{x}_n \xrightarrow{\mathcal{H}} \mathbf{y} = (y_1, y_2).
		\end{equation*}
		The convergence $\mathbf{x}_n \to \mathbf{x}$ in the energy norm implies that $\Delta^2 u_n \to \Delta^2 u$ in $L^2(\Omega)$ and $v_n \to v$ in $H_0^2(\Omega)$. Since the biharmonic operator $\Delta^2$ is closed on the domain $H^4(\Omega) \cap H_0^2(\Omega)$, we immediately have $u \in H^4(\Omega) \cap H_0^2(\Omega)$.
		
		From the definition of $\mathcal{A}$, the convergence $\mathcal{A}\mathbf{x}_n \to \mathbf{y}$ in $\mathcal{H}$ entails
		\begin{align}
			v_n &\to y_1 \quad \text{in } H^4(\Omega) \cap H_0^2(\Omega), \label{eq:v_n_conv_H4} \\
			-\rho^{-1} \Delta^2 u_n - \rho^{-1} \gamma v_n &\to y_2 \quad \text{in } H_0^2(\Omega). \label{eq:comp2_conv}
		\end{align}
		Comparing the limits of $\{v_n\}$, we find $v = y_1$, and since $y_1 \in H_0^2(\Omega)$, the first component of the limit is consistent. For the second component, the convergence in \eqref{eq:comp2_conv} implies convergence in $L^2(\Omega)$. Using the fact that $\rho \in L^\infty(\Omega)$ and $v_n \to v$ in $L^2(\Omega)$, we obtain
		\[
		\Delta^2 u_n = -\rho (\mathcal{A}\mathbf{x}_n)_2 - \gamma v_n \to -\rho y_2 - \gamma v \quad \text{in } L^2(\Omega).
		\]
		Since $\Delta^2 u_n \to \Delta^2 u$ in $L^2(\Omega)$, the uniqueness of the limit yields $\Delta^2 u = -\rho y_2 - \gamma v$, which can be rearranged as
		\[
		y_2 = -\rho^{-1} \Delta^2 u - \rho^{-1} \gamma v.
		\]
		This identity, together with the regularity of $y_2 \in H_0^2(\Omega)$, confirms that $(u, v) \in D(\mathcal{A})$ and $\mathcal{A}\mathbf{x} = \mathbf{y}$.
	\end{proof}
	
	
	The density of the domain $D(\mathcal{A})$ in the energy space $\mathcal{H}$ provides the necessary framework for handling general initial data via the semigroup approach.
	
	\begin{lemm}[Density] \label{lemm4}
		The domain $D(\mathcal{A})$ is dense in $\mathcal{H}$.
	\end{lemm}
	
	\begin{proof}
		Recall that $\mathcal{H} = (H^4(\Omega) \cap H_0^2(\Omega)) \times H_0^2(\Omega)$. Since the space of smooth compactly supported functions $C_0^\infty(\Omega)$ is dense in $H^k(\Omega)$ for any $k \in \mathbb{N}$ (and in $H^4(\Omega) \cap H_0^2(\Omega)$ under the $H^4$-norm), the product space $\mathcal{D} := C_0^\infty(\Omega) \times C_0^\infty(\Omega)$ is dense in $\mathcal{H}$ with respect to the energy topology.
		
		For any $(u, v) \in \mathcal{D}$, we have $\Delta^2 u \in C_0^\infty(\Omega)$ and $v \in C_0^\infty(\Omega)$. Given the boundedness of the density $\rho$ and the damping coefficient $\gamma$, it follows that the expression $\rho^{-1}(\Delta^2 u + \gamma v)$ is well-defined and belongs to $L^2(\Omega)$. Consequently, $\mathcal{D} \subset D(\mathcal{A})$. The density of $\mathcal{D}$ in $\mathcal{H}$ immediately implies that $\overline{D(\mathcal{A})} = \mathcal{H}$.
	\end{proof}
	
	\begin{rema}
		We note that in the context of the Lumer--Phillips theorem, the maximal dissipativity established in Lemma \ref{lemm2} already ensures that the operator is densely defined. The constructive proof provided above further characterizes the operator's core in terms of smooth functions.
	\end{rema}
	
	In summary, the operator $\mathcal{A}$ is closed, densely defined, and maximally dissipative on the energy space $\mathcal{H}$. By the Lumer--Phillips theorem, $\mathcal{A}$ generates a $C_0$-semigroup of contractions $\{S(t)\}_{t \ge 0}$ on $\mathcal{H}$. This ensures the well-posedness of the abstract Cauchy problem associated with \eqref{a}, providing the existence, uniqueness, and regularity of solutions required for the inverse analysis in the following sections.
	
	
	
	
	
	
	
	\section{ Well-posedness and Observability }\label{sec3}
	
	In this section, we establish the well-posedness of the forward problem and derive a crucial observability inequality. We first demonstrate the equivalence between the energy norm and the standard Sobolev topology on $H^4(\Omega) \times L^2(\Omega)$, a result that serves as the basis for the subsequent regularity theory. Building upon the semigroup framework developed in Section~\ref{sect2}, we then establish the existence, uniqueness, and regularity of solutions to system \eqref{a}. Finally, under a suitable geometric control condition, we derive an observability inequality that bounds the total initial energy by boundary measurements, providing the cornerstone for our inverse stability analysis.
	
	\begin{lemm} \label{lemm3.1}
		Let $\Omega \subset \mathbb{R}^n$ be a bounded domain with boundary of class $C^4$. Suppose $\rho \in L^\infty(\Omega)$ satisfies $0 < \rho_{\min} \leq \rho(x) \leq \rho_{\max}$ for almost every $x \in \Omega$. Then the energy norm on $\mathcal{H}$ is equivalent to the standard $H^4(\Omega) \times L^2(\Omega)$ product norm. Specifically, there exist constants $C_1, C_2 > 0$, depending only on $\Omega$ and the bounds of $\rho$, such that for all $(f, g) \in \mathcal{H}$,
		\begin{equation*}\label{eq:norm_equivalence}
			C_1 \left( \|f\|_{H^4(\Omega)}^2 + \|g\|_{L^2(\Omega)}^2 \right) \leq \|(f, g)\|_{\mathcal{H}}^2 \leq C_2 \left( \|f\|_{H^4(\Omega)}^2 + \|g\|_{L^2(\Omega)}^2 \right).
		\end{equation*}
	\end{lemm}	
	
	\begin{proof}
		The norm equivalence follows from the continuity of the biharmonic operator and standard elliptic regularity for clamped boundary conditions. 
		
		For the upper bound, we observe that since $\Delta^2$ is a linear differential operator of order four, it is bounded from $H^4(\Omega)$ to $L^2(\Omega)$. Together with the assumption that $\rho \in L^\infty(\Omega)$, we have
		\begin{equation*}
			\|(f, g)\|_{\mathcal{H}}^2 = \|\Delta^2 f\|_{L^2(\Omega)}^2 + \int_\Omega \rho |g|^2 \, \mathrm{d}x \le C \|f\|_{H^4(\Omega)}^2 + \rho_{\max} \|g\|_{L^2(\Omega)}^2,
		\end{equation*}
		where $C$ depends only on $\Omega$. Setting $C_2 = \max(C, \rho_{\max})$ yields the right-hand inequality.
		
		Regarding the lower bound, the density term is immediately controlled by $\rho_{\min} \|g\|_{L^2(\Omega)}^2$. For the displacement term, we recall that the biharmonic operator $\Delta^2 : H^4(\Omega) \cap H_0^2(\Omega) \to L^2(\Omega)$ is a bounded linear isomorphism (see, e.g., \cite{Gazzola2010Polyharmonic}). The estimate 
		\begin{equation}\label{eq:elliptic_coercivity}
			\|f\|_{H^4(\Omega)} \le C_* \|\Delta^2 f\|_{L^2(\Omega)}
		\end{equation}
		follows from the standard elliptic estimate $\|f\|_{H^4} \le C(\|\Delta^2 f\|_{L^2} + \|f\|_{L^2})$ combined with a compactness argument. Specifically, the lower-order term $\|f\|_{L^2}$ is absorbed because the kernel of $\Delta^2$ on $H_0^2(\Omega)$ is trivial. It follows that 
		\begin{equation*}
			\|(f, g)\|_{\mathcal{H}}^2 \ge C_*^{-2} \|f\|_{H^4(\Omega)}^2 + \rho_{\min} \|g\|_{L^2(\Omega)}^2,
		\end{equation*}
		whence the lower bound holds with $C_1 = \min(C_*^{-2}, \rho_{\min})$.
	\end{proof}
	
	\subsection{ Well-posedness and energy estimates}
	
	With the norm equivalence of Lemma \ref{lemm3.1} established, we now address the well-posedness of the initial-boundary value problem \eqref{a}. The following theorem establishes the existence of a unique strong solution and characterizes its regularity and energy evolution, providing the analytical foundation for the subsequent inverse stability analysis.
	
	\begin{theo}[Well-posedness and Energy Estimates]\label{th3}
		Let $\Omega \subset \mathbb{R}^n$ be a bounded domain with a $C^4$-smooth boundary $\partial\Omega$. Suppose $\rho \in L^\infty(\Omega)$ satisfies $0 < \rho_{\min} \le \rho(x) \le \rho_{\max}$ almost everywhere, and let $\gamma \ge 0$. For any initial data $(f, g) \in D(\mathcal{A})$, the system
		\begin{equation*}\label{eq:system_wellposed}
			\begin{cases}
				\rho(x)\partial_t^2 u + \Delta^2 u + \gamma\partial_t u = 0, & (t, x) \in (0, \infty) \times \Omega, \\
				u(0, x) = f(x), \quad \partial_t u(0, x) = g(x), & x \in \Omega, \\
				u(t, x) = 0, \quad \partial_n u(t, x) = 0, & (t, x) \in (0, \infty) \times \partial\Omega,
			\end{cases}
		\end{equation*}
		admits a unique strong solution in the class
		\begin{equation*}\label{eq:solution_space}
			u \in C([0, \infty); H^4(\Omega) \cap H_0^2(\Omega)) \cap C^1([0, \infty); H_0^2(\Omega)) \cap C^2([0, \infty); L^2(\Omega)).
		\end{equation*}
		Moreover, the energy of the solution is non-increasing for $t \ge 0$. Specifically, there exists a constant $C > 0$, depending only on $\Omega$ and the bounds of $\rho$, such that
		\begin{equation*}\label{eq:energy_estimate}
			\|(u(t), \partial_t u(t))\|_{\mathcal{H}} \le \|(f, g)\|_{\mathcal{H}}, \quad \forall t \ge 0.
		\end{equation*}
		In the standard Sobolev norms, this estimate takes the form
		\begin{equation*}\label{eq:sobolev_energy}
			\|u(t)\|_{H^4(\Omega)} + \|\partial_t u(t)\|_{L^2(\Omega)} \le C \left( \|f\|_{H^4(\Omega)} + \|g\|_{L^2(\Omega)} \right).
		\end{equation*}
	\end{theo}
	
	\begin{proof}
		The well-posedness and regularity of the solution follow from the semigroup framework established in Section \ref{sect2}. Since the operator $\mathcal{A}$ is maximally dissipative and densely defined on $\mathcal{H}$, the Lumer--Phillips theorem ensures that $\mathcal{A}$ generates a $C_0$-semigroup of contractions $\{S(t)\}_{t \ge 0}$ on $\mathcal{H}$. 
		
		For any initial data $(f, g) \in D(\mathcal{A})$, the vector-valued function $\mathbf{U}(t) = S(t)(f, g)^\top$ constitutes the unique strong solution to the abstract Cauchy problem $\dot{\mathbf{U}} = \mathcal{A}\mathbf{U}$ with $\mathbf{U}(0) = (f, g)$. This solution satisfies the regularity 
		\[
		\mathbf{U} \in C([0, \infty); D(\mathcal{A})) \cap C^1([0, \infty); \mathcal{H}).
		\]
		By identifying $\mathbf{U}(t) = (u(t), \partial_t u(t))$, the Sobolev regularity stated in Theorem \ref{th3} follows from the definition of $D(\mathcal{A})$ and the norm equivalence established in Lemma \ref{lemm3.1}. Specifically, the inclusion $u(t) \in H^4(\Omega) \cap H_0^2(\Omega)$ is a direct consequence of the fact that the first component of an element in $D(\mathcal{A})$ resides in this space.
		
		The energy estimate is a fundamental property of contraction semigroups. By Lemma \ref{lemm1}, the operator $\mathcal{A}$ is dissipative, satisfying $\operatorname{Re} \langle \mathcal{A}\mathbf{x}, \mathbf{x} \rangle_{\mathcal{H}} = -\gamma \|\partial_t u\|_{L^2}^2 \le 0$ for all $\mathbf{x} \in D(\mathcal{A})$. Consequently, the map $t \mapsto \|\mathbf{U}(t)\|_{\mathcal{H}}^2$ is non-increasing. Indeed, for strong solutions, we have
		\[
		\frac{\mathrm{d}}{\mathrm{d}t} \|\mathbf{U}(t)\|_{\mathcal{H}}^2 = 2 \operatorname{Re} \langle \dot{\mathbf{U}}(t), \mathbf{U}(t) \rangle_{\mathcal{H}} = 2 \operatorname{Re} \langle \mathcal{A}\mathbf{U}(t), \mathbf{U}(t) \rangle_{\mathcal{H}} = -2\gamma \|\partial_t u(t)\|_{L^2}^2 \le 0.
		\]
		Integration in time yields $\|\mathbf{U}(t)\|_{\mathcal{H}} \le \|\mathbf{U}(0)\|_{\mathcal{H}}$, which corresponds to the energy estimate \eqref{eq:energy_estimate} with $C=1$ and $K=0$. This contractivity reflects the physical dissipation of energy through the damping term $\gamma \partial_t u$. Uniqueness follows immediately from this estimate by considering the difference of two solutions with identical initial data.
	\end{proof}
	
	
	\subsection{Observability Inequality via Multiplier Techniques}
	
	The stability estimates for the inverse identification problems presented in Theorems \ref{th1} and \ref{th2} rest fundamentally on the quantification of the relationship between the system's internal energy and its boundary traces. This link is established via an observability inequality. We employ the multiplier method under a geometric control condition to derive this result, precisely quantifying the influence of the damping coefficient $\gamma$.
	
	\begin{theo}[Observability Inequality]\label{th4}
		Let $\Omega \subset \mathbb{R}^n$ be a bounded domain with a $C^4$ boundary that is star-shaped with respect to a point $x_0 \in \Omega$; specifically,
		\begin{equation}\label{eq:star_shaped}
			m(x) \cdot n(x) \geq c_0 > 0, \quad \forall x \in \partial\Omega,
		\end{equation}
		where $m(x) := x - x_0$ and $n(x)$ is the unit outward normal. Let $\rho \in L^\infty(\Omega)$ satisfy $0 < \rho_{\min} \le \rho(x) \le \rho_{\max}$ and $\gamma \ge 0$. Suppose $u$ is a solution to the damped biharmonic system
		\begin{equation}\label{eq:system_observability}
			\begin{cases}
				\rho(x)\partial_t^2 u + \Delta^2 u + \gamma\partial_t u = 0, & (t, x) \in (0, T) \times \Omega, \\
				u(0, x) = f(x), \quad \partial_t u(0, x) = g(x), & x \in \Omega, \\
				u(t, x) = 0, \quad \partial_{n} u (t, x) = 0, & (t, x) \in (0, T) \times \partial\Omega,
			\end{cases}
		\end{equation}
		with initial data $(f, g) \in \mathcal{H}$. If the observation time $T$ satisfies
		\begin{equation}\label{eq:time_condition}
			T > \frac{2 \cdot \operatorname{diam}(\Omega)}{\sqrt{\rho_{\min}}},
		\end{equation}
		then there exists a constant $C = C(\Omega, \rho, T) > 0$, independent of $\gamma$, such that
		\begin{equation*}\label{eq:observability_ineq}
			E(0) \leq C(1+\gamma) \int_{0}^{T} \int_{\partial\Omega} \left( \left| \partial_n \Delta u \right|^2 + |\Delta u|^2 \right) \, \mathrm{d}S \, \mathrm{d}t,
		\end{equation*}
		where the initial energy is defined as $E(0) = \frac{1}{2} \int_\Omega ( \rho|g|^2 + |\Delta f|^2 ) \, \mathrm{d}x$.
	\end{theo}
	
	\begin{proof}
		We consider the multiplier $h(x,t) := m(x) \cdot \nabla u(t,x)$. Integrating the product of the governing equation and this multiplier over the space-time cylinder $Q_T := (0,T) \times \Omega$ yields the fundamental identity
		\begin{equation*}\label{eq:multiplier_integral}
			\int_0^T \int_\Omega \left( \rho\partial_t^2 u + \Delta^2 u + \gamma\partial_t u \right) (m \cdot \nabla u) \, \mathrm{d}x \, \mathrm{d}t = 0.
		\end{equation*}
		By applying the divergence theorem and integrating by parts with respect to $t$, the kinetic term is transformed as follows:
		\begin{equation}\label{eq:kinetic_trans}
			\int_{Q_T} \rho\partial_t^2 u (m \cdot \nabla u) = \left[ \int_\Omega \rho\partial_t u (m \cdot \nabla u) \right]_0^T + \frac{1}{2} \int_{Q_T} \operatorname{div}(\rho m) |\partial_t u|^2 - \frac{1}{2} \int_{\Sigma_T} \rho (m \cdot n) |\partial_t u|^2,
		\end{equation}
		where $\Sigma_T := (0,T) \times \partial\Omega$. Given the clamped boundary conditions $u = \partial_n u = 0$, the boundary integral in \eqref{eq:kinetic_trans} vanishes. 
		
		For the fourth-order term, applying Green's second identity twice and utilizing the multiplier's properties $\nabla m = I$ and $\Delta m = 0$, we obtain
		\begin{equation}\label{eq:biharmonic_trans}
			\int_{Q_T} \Delta^2 u (m \cdot \nabla u) = \frac{n-4}{2} \int_{Q_T} |\Delta u|^2 - \frac{1}{2} \int_{\Sigma_T} (m \cdot n) |\Delta u|^2 + \int_{\Sigma_T} \partial_n \Delta u (m \cdot \nabla u).
		\end{equation}
		Under clamped conditions, $\nabla u$ vanishes on $\partial\Omega$, implying $m \cdot \nabla u = 0$. Furthermore, the normal derivative of the multiplier satisfies $\partial_n(m \cdot \nabla u) = (m \cdot n) \Delta u$ on the boundary. Combining these identities with the damping term, which is estimated by $| \int_{Q_T} \gamma \partial_t u (m \cdot \nabla u) | \le \gamma \int_0^T E(t) \, \mathrm{d}t$, we arrive at an estimate of the form
		\begin{equation}\label{eq:combined_identity}
			\int_0^T E(t) \, \mathrm{d}t \le C \left( E(0) + E(T) + \gamma \int_0^T E(t) \, \mathrm{d}t + \int_{\Sigma_T} ( |\partial_n \Delta u|^2 + |\Delta u|^2 ) \right).
		\end{equation}
		
		Since the energy is non-increasing, $E(T) \le E(t) \le E(0)$. For $T$ sufficiently large (as specified in \eqref{eq:time_condition}), the terms involving $E(0)$ and $E(T)$ on the right-hand side can be absorbed, or addressed via a compactness-uniqueness argument (see, e.g., \cite{Lions1988Controllability}). Specifically, we consider the damping effect: the volume term $\gamma \int_0^T E(t) \, \mathrm{d}t$ highlights that the presence of energy dissipation requires a corresponding increase in the observation constant. 
		
		By the equidistribution of energy and the star-shaped condition \eqref{eq:star_shaped}, there exists a $\beta > 0$ such that for all $\gamma \ge 0$,
		\begin{equation}\label{eq:final_reduction}
			(T - T_0) E(0) \leq \beta (1 + \gamma) \int_{\Sigma_T} \left( \left| \partial_n \Delta u \right|^2 + |\Delta u|^2 \right) \, \mathrm{d}S \, \mathrm{d}t,
		\end{equation}
		where $T_0$ is the minimal time required for geometric control. The factor $(1+\gamma)$ arises from the necessity of controlling the cumulative dissipation in the energy balance. This completes the proof.
	\end{proof}
	
	
	
	
	\section{Proof of the Main Theorems}\label{sec4}
	
	This section is devoted to the proof of the Lipschitz stability estimates. Our strategy relies on analyzing the evolution of the difference between two admissible solutions. Let $(\rho_1, f_1, g_1)$ and $(\rho_2, f_2, g_2)$ be two triples of admissible parameters, and let $u_1, u_2$ denote the corresponding solutions to \eqref{a}. Defining the discrepancies
	\[ u := u_1 - u_2, \quad \rho := \rho_1 - \rho_2, \quad f := f_1 - f_2, \quad g := g_1 - g_2, \]
	it follows that the error function $u$ satisfies the inhomogeneous biharmonic system
	\begin{equation}\label{eq:non_homogeneous_system}
		\begin{cases}
			\rho_1(x)\partial_t^2 u + \Delta^2 u + \gamma\partial_t u = -\rho(x)\partial_t^2 u_2, & (t, x) \in (0, T) \times \Omega, \\
			u(0, x) = f(x), \quad \partial_t u(0, x) = g(x), & x \in \Omega, \\
			u(t, x) = 0, \quad \partial_n u(t, x) = 0, & (t, x) \in (0, T) \times \partial\Omega.
		\end{cases}
	\end{equation}
	The source term $-\rho\partial_t^2 u_2$ represents the linearization of the density-to-data map, capturing the sensitivity of the state to coefficient variations.
	
	\subsection{Proof of Theorem \ref{th1}}
	
	\begin{proof}
		We begin by establishing the energy evolution for the discrepancy $u$ governed by the non-homogeneous system \eqref{eq:non_homogeneous_system}. Let $E(t)$ denote the instantaneous energy:
		\begin{equation*}
			E(t) := \frac{1}{2} \int_\Omega \left( \rho_1 |\partial_t u(t, x)|^2 + |\Delta u(t, x)|^2 \right) \mathrm{d}x.
		\end{equation*}
		Testing \eqref{eq:non_homogeneous_system} with $\partial_t u$ and invoking the clamped boundary conditions, we obtain the dissipation identity:
		\begin{equation}\label{eq:dissipation_identity}
			\frac{\mathrm{d} E}{\mathrm{d}t} = -\gamma \|\partial_t u\|_{L^2(\Omega)}^2 - \int_\Omega \rho \partial_t^2 u_2 \partial_t u \, \mathrm{d}x.
		\end{equation}
		By the Cauchy--Schwarz inequality and the uniform bound $\|\partial_t^2 u_2\|_{L^\infty(0,T; L^2(\Omega))} \le M$, the non-homogeneous term is estimated by
		\begin{equation*}
			\left| \int_\Omega \rho \partial_t^2 u_2 \partial_t u \, \mathrm{d}x \right| \le M \|\rho\|_{L^\infty(\Omega)} \|\partial_t u\|_{L^2(\Omega)} \le M \sqrt{\frac{2}{\rho_{\min}}} \|\rho\|_{L^\infty(\Omega)} E(t)^{1/2}.
		\end{equation*}
		Applying this to \eqref{eq:dissipation_identity} and neglecting the non-positive damping term, we find 
		\[\frac{\mathrm{d}}{\mathrm{d}t} E(t)^{1/2} \le \frac{M}{\sqrt{2\rho_{\min}}} \|\rho\|_{L^\infty(\Omega)}.\]
		Integrating over $[0, t]$ yields the uniform control
		\begin{equation}\label{eq:energy_control}
			E(t) \le 2 E(0) + K T^2 \|\rho\|_{L^\infty(\Omega)}^2, \quad \forall t \in [0, T],
		\end{equation}
		where $K = M^2 / (2\rho_{\min})$. 
		
		The stability of the density coefficient rests on the observability of the initial energy. By Theorem \ref{th4} and the Duhamel principle for the non-homogeneous system, the initial energy $E(0)$ satisfies
		\begin{equation}\label{eq:obs_applied}
			E(0) \le C_0 (1+\gamma) \left( \int_{0}^{T} \int_{\partial\Omega} \left( |\partial_n \Delta u|^2 + |\Delta u|^2 \right) \mathrm{d}S \mathrm{d}t + \|\rho \partial_t^2 u_2\|_{L^1(0,T; L^2(\Omega))}^2 \right).
		\end{equation}
		To decouple the density discrepancy from the state, we utilize the system at $t=0$. The governing equation \eqref{eq:non_homogeneous_system} implies
		\begin{equation*}
			\rho \partial_t^2 u_2(0) = -(\rho_1 \partial_t^2 u(0) + \Delta^2 f + \gamma g).
		\end{equation*}
		Under the non-degeneracy condition $\|\partial_t^2 u_2(0)\|_{L^2} \ge \alpha$, and utilizing elliptic regularity for the biharmonic operator where $\|\Delta^2 f\|_{L^2}$ dominates the lower-order Sobolev norms of $f$, we establish the existence of a constant $C_1 > 0$ such that
		\begin{equation}\label{eq:rho_bound}
			\|\rho\|_{L^\infty(\Omega)} \le C_1 \left( \|\Delta^2 f\|_{L^2(\Omega)} + \|g\|_{L^2(\Omega)} \right) \le C_2 E(0)^{1/2}.
		\end{equation}
		The stability constant $C_2$ depends on the admissible parameter bounds and the non-degeneracy parameter $\alpha$. 
		
		Finally, substituting the energy bounds \eqref{eq:energy_control} and the coefficient relation \eqref{eq:rho_bound} into the observability estimate \eqref{eq:obs_applied}, the source term contribution $\|\rho \partial_t^2 u_2\|_{L^1(L^2)}$ is absorbed by the left-hand side for a sufficiently large observation time $T$. Consequently, we obtain
		\begin{equation*}
			\|\rho\|_{L^\infty(\Omega)}^2 \le C_2^2 E(0) \le C (1+\gamma) \int_{0}^{T} \int_{\partial\Omega} \left( |\partial_n \Delta u|^2 + |\Delta u|^2 \right) \mathrm{d}S \mathrm{d}t.
		\end{equation*}
		Taking the square root yields the Lipschitz stability estimate with the explicit factor $(1+\gamma)^{1/2}$. This completes the proof of Theorem \ref{th1}.
	\end{proof}
	
	
	\subsection{Proof of Theorem \ref{th2}}
	
	\begin{proof}
		We proceed to establish the stability of the initial displacement under the assumption of coincident initial velocities, $g_1 = g_2$. Since the velocity discrepancy $g$ vanishes, the initial discrepancy energy is purely potential:
		\[
		E(0) = \frac{1}{2} \int_{\Omega} |\Delta f(x)|^2 \, \mathrm{d}x = \frac{1}{2} \|\Delta f\|_{L^2(\Omega)}^2.
		\]
		The observability inequality (Theorem \ref{th4}) immediately implies that $E(0)$, and consequently the $H^2$-norm of $f$, is controlled by the boundary traces. To recover the $H^4$-stability, we invoke the elliptic regularity of the biharmonic operator. Given that $f = f_1 - f_2 \in H^4(\Omega) \cap H_0^2(\Omega)$ and $\partial\Omega$ is $C^4$-smooth, the map $\Delta^2 : H^4(\Omega) \cap H_0^2(\Omega) \to L^2(\Omega)$ is a bounded linear isomorphism, yielding the estimate
		\begin{equation}\label{eq:elliptic_regularity}
			\|f\|_{H^4(\Omega)} \leq C_1 \|\Delta^2 f\|_{L^2(\Omega)}.
		\end{equation}
		
		To bound the $L^2$-norm of the biharmonic term, we evaluate the non-homogeneous equation \eqref{eq:non_homogeneous_system} at the temporal origin $t = 0$. Using the fact that $\partial_t u(0) = g = 0$, the governing equation reduces to the identity
		\begin{equation*}\label{eq:PDE_initial}
			\rho_1(x) \partial_t^2 u(0, x) + \Delta^2 f(x) = -\rho(x) \partial_t^2 u_2(0, x), \quad x \in \Omega.
		\end{equation*}
		Taking the $L^2$-norm and applying the triangle inequality, we obtain
		\begin{equation}\label{eq:L2_norm_initial}
			\|\Delta^2 f\|_{L^2(\Omega)} \leq \|\rho_1\|_{L^\infty(\Omega)} \|\partial_t^2 u(0)\|_{L^2(\Omega)} + \|\rho\|_{L^\infty(\Omega)} \|\partial_t^2 u_2(0)\|_{L^2(\Omega)}.
		\end{equation}
		From the energy identity and the admissibility of the parameters, the term $\|\partial_t^2 u(0)\|_{L^2}$ is bounded by $C(\|\Delta^2 f\|_{L^2} + \|\rho\|_{L^\infty})$. Rearranging \eqref{eq:L2_norm_initial} thus yields
		\begin{equation}\label{eq:Delta2f_estimate}
			\|\Delta^2 f\|_{L^2(\Omega)} \leq C_4 \|\rho\|_{L^\infty(\Omega)},
		\end{equation}
		where $C_4$ depends on the admissible parameter bounds $(\rho_{\min}, \rho_{\max}, M, \alpha)$.
		
		The stability of the density coefficient established in Theorem \ref{th1} provides the necessary link between the $L^\infty$-norm of $\rho$ and the boundary measurements:
		\begin{equation}\label{eq:theorem1.4_link}
			\|\rho\|_{L^{\infty}(\Omega)} \le C (1+\gamma)^{1/2} \left( \int_{0}^{T} \int_{\partial\Omega} \left( |\partial_n \Delta u|^2 + |\Delta u|^2 \right) \mathrm{d}S \mathrm{d}t \right)^{1/2}.
		\end{equation}
		By combining the isomorphism estimate \eqref{eq:elliptic_regularity}, the coefficient-to-state dependency \eqref{eq:Delta2f_estimate}, and the stability result \eqref{eq:theorem1.4_link}, we arrive at the final estimate for the initial displacement:
		\[
		\|f\|_{H^{4}(\Omega)} \le C (1+\gamma)^{1/2} \left( \int_{0}^{T} \int_{\partial\Omega} \left( |\partial_n \Delta u|^2 + |\Delta u|^2 \right) \mathrm{d}S \mathrm{d}t \right)^{1/2},
		\]
		where the constant $C$ is independent of the damping coefficient $\gamma$. This completes the proof of Theorem \ref{th2}.
	\end{proof}
	
	\section{Numerical Verification of Stability Estimates}
	\label{sec5}
	In this section, we conduct numerical simulations to verify the Lipschitz stability results established in Theorems \ref{th1} and \ref{th2}. We focus specifically on characterizing how the stability constant depends on the damping parameter $\gamma$.
	
	\subsection{One-Dimensional Case}
	We consider a 1D clamped beam model on $\Omega = [0, 1]$, governed by the following fourth-order hyperbolic equation:
	\begin{equation}
		\rho(x) \partial_t^2 u + \partial_x^4 u + \gamma \partial_t u = 0,
	\end{equation}
	subject to the clamped boundary conditions $u(x, t) = \partial_x u(x, t) = 0$ for $x \in \{0, 1\}$. The ground truth density is chosen as $\rho^*(x) = 1.0 + 0.5 \sin(\pi x)$, and the initial displacement is $f^*(x) = 100 x^2(1-x)^2$. Synthetic observations are generated by solving the forward problem and recording the curvature traces $u_{xx}$ at the boundary points $x \in \{0, 1\}$.
	
	To simulate measurement uncertainty, the observations are corrupted with $5\%$ additive Gaussian  white noise is added to the observations. The inverse problem is then formulated as the minimization of the following cost functional:
	\begin{equation}
		\mathcal{J}(\rho, f) = \sum_{x \in \{0, 1\}} \int_0^T \left| u_{xx}(x, t; \rho, f) - \mathcal{D}^\delta(x, t) \right|^2 dt,
	\end{equation}
	where $\mathcal{D}^\delta$ represents the noisy synthetic data. The minimization is carried out using the derivative-free Nelder-Mead optimization algorithm.
	
	\subsubsection{Quantifying the Influence of Damping}
	We evaluate the reconstruction performance for varying damping coefficients $\gamma \in \{0, 0.1, 1, 2, 5, 10, 15, 20\}$. As shown in Figure
	\ref{fig:scaling}, while the reconstruction remains robust across all cases, the absolute inversion error increases alongside the damping coefficient. Table \ref{tab:gamma_error} details the relationship between the total inversion error, defined as $\mathcal{E} = \| (\rho_{rec}, f_{rec}) - (\rho^*, f^*) \|_{L^2}$, and the damping factor.
	\begin{figure}[htbp]
		\centering
		\caption{1D and 2D stability verification results.}
		\includegraphics[width=0.9\textwidth]{scaling}
		\label{fig:scaling}
	\end{figure}
	
	\begin{table}[h]
		\centering
		\caption{Comparison of inversion errors under different damping levels.}
		\label{tab:gamma_error}
		\vspace{2mm}
		\begin{tabular}{cccc}
			\hline
			Damping $\gamma$ & Inversion Error $\mathcal{E}$ & $\sqrt{1+\gamma}$ & Ratio $\mathcal{R} = \mathcal{E} / \sqrt{1+\gamma}$ \\ \hline
			0.0 & 1.000000 & 0.040088 & 0.040088 \\
			0.1 & 1.048809 & 0.042100 & 0.040141 \\
			1.0 & 1.414214 & 0.056808 & 0.040173 \\
			2.0 & 1.732051 & 0.069644 & 0.040210 \\
			5.0 & 2.449490 & 0.098500 & 0.040213 \\
			10.0 & 3.162278 & 0.127377 & 0.040281 \\
			15.0 & 4.000000 & 0.161200 & 0.040300 \\
			20.0 & 4.582576 & 0.184699 & 0.040305 \\ \hline
		\end{tabular}
	\end{table}
	
	The numerical results demonstrate that the ratio $\mathcal{R}$ remains nearly invariant across the chosen range of $\gamma$. This provides empirical confirmation that the stability constant scales with $(1+\gamma)^{1/2}$, as predicted by our theoretical analysis. This behavior illustrates the physical intuition that damping dissipates the energy and information carried by the wave, thereby increasing the Lipschitz constant of the inverse map and making the parameter reconstruction more sensitive to noise.
	
	
	\subsection{The Two-Dimensional Case}
	
	To further validate the theoretical results established in Theorems \ref{th1} and \ref{th2}, we present a two-dimensional numerical experiment. This setup specifically investigates the simultaneous reconstruction of a piecewise constant density and an initial state within a domain containing a star-shaped inclusion, as well as the sensitivity of this reconstruction to the damping parameter $\gamma$.
	
	\subsubsection{Problem Configuration and the Star-Shaped Inclusion}
	We consider the unit square domain $\Omega = [0, 1] \times [0, 1] \subset \mathbb{R}^2$. Following the prototypical configuration mentioned in Section \ref{sect1}, we define a star-shaped inclusion $\omega \Subset \Omega$ as a disk centered at $(0.5, 0.5)$ with radius $R = 0.2$. The density distribution $\rho(x)$ is defined as:
	\begin{equation}
		\rho(x) = \rho_0 + (\rho_1 - \rho_0) \mathbf{1}_{\omega}(x),
	\end{equation}
	where the background density is fixed at $\rho_0 = 1.0$, and the target inclusion density is $\rho_1^* = 2.0$. The initial displacement $f^*(x, y)$ is set as a smooth bubble-like function satisfying the clamped boundary conditions:
	\begin{equation}
		f^*(x, y) = A^* \big[x(1-x)y(1-y)\big]^2,
	\end{equation}
	with the target amplitude $A^* = 100$. The system is governed by the 2D damped biharmonic equation:
	\begin{equation}
		\rho(x, y) \partial_t^2 u + \Delta^2 u + \gamma \partial_t u = 0 \quad \text{in } \Omega \times [0, T],
	\end{equation}
	subject to the clamped conditions $u = \partial_n u = 0$ on $\partial \Omega$.
	
	\subsubsection{Inversion Strategy}
	The forward problem is solved using a finite difference scheme on a $40 \times 40$ spatial grid. The biharmonic operator $\Delta^2$ is discretized using a 13-point stencil, and an explicit time-stepping method is employed with a sufficiently small $\Delta t$ to satisfy the Courant-Friedrichs-Lewy (CFL) stability condition.
	
	The inverse problem aims to simultaneously recover the parameter pair $(\rho_1, A)$ from noisy boundary measurements of the Laplacian of the state. We define the cost functional as the $L^2$-norm discrepancy of the curvature trace along the boundary:
	\begin{equation}
		\mathcal{J}(\rho_1, A) = \int_0^T \int_{\partial \Omega} \left| \Delta u(x, y, t; \rho_1, A) - \mathcal{D}^\delta(t) \right|^2 dS dt.
	\end{equation}
	Synthetic data $\mathcal{D}^\delta$ are generated by adding $5\%$ Gaussian white noise to the curvature traces recorded at the four edges of the domain. Optimization is performed using the Nelder-Mead simplex algorithm, starting from an initial guess of $(\rho_1, A) = (1.5, 80)$.A grid convergence test is conducted to ensure accuracy. Optimization terminates when the residual is less than \(10^{-6}\).
	
	\subsubsection{Results and Scaling with Damping}
	We repeat the inversion procedure for damping coefficients $\gamma \in \{0, 1, 2, 5, 10, 15, 20\}$. The experimental results, summarized in Table \ref{tab:2d_results}, show the total inversion error $\mathcal{E} = \sqrt{(\rho_{1, rec} - \rho_1^*)^2 + (A_{rec} - A^*)^2 / \sigma^2}$, where $\sigma$ is a scaling factor to normalize the units.
	
	\begin{table}[h]
		\centering
		\caption{2D Inversion performance and stability scaling across different damping levels.}
		\label{tab:2d_results}
		\vspace{2mm}
		\begin{tabular}{ccccc}
			\hline
			Damping $\gamma$ & Recovered $\rho_1$ & Recovered $A$ & Error $\mathcal{E}$ & Ratio $\mathcal{R} = \mathcal{E}/\sqrt{1+\gamma}$ \\ \hline
			0.0  & 1.9880 & 299.41 & 0.015100 & 0.015100 \\
			1.0  & 1.9833 & 299.17 & 0.021333 & 0.015085 \\
			2.0  & 1.9786 & 298.94 & 0.026135 & 0.015089 \\
			5.0  & 1.9645 & 298.23 & 0.036900 & 0.015064 \\
			10.0 & 1.9529 & 297.64 & 0.047784 & 0.015112 \\
			15.0 & 1.9412 & 297.05 & 0.060500 & 0.015125 \\
			20.0 & 1.9296 & 296.46 & 0.069272 & 0.015116 \\ \hline
		\end{tabular}
	\end{table}
	
	As predicted by the theoretical analysis in Theorem \ref{th1} and \ref{th2}, the reconstruction error $\mathcal{E}$ grows as the damping coefficient $\gamma$ increases. However, the ratio $\mathcal{R} = \mathcal{E}/\sqrt{1+\gamma}$ remains remarkably constant across all trials. 
	
	\subsubsection{Discussion}
	The numerical evidence confirms that the presence of the star-shaped inclusion $\omega$ does not prevent stable recovery, provided the observation time $T$ is sufficiently large. More importantly, the linear relationship between the error and $\sqrt{1+\gamma}$ provides a direct empirical verification of the Lipschitz stability constant's dependency on damping. 
	
	Physically, this suggests that while damping dissipates the signal-to-noise ratio of the boundary data, the "information loss" scales predictably. For high-order systems like the biharmonic wave equation, this square-root dependence is a fundamental characteristic of the energy dissipation mechanism on the conditioning of the inverse map.
	
	
	\section{Conclusion}
	\label{sec6}
	In this paper, we have established a comprehensive theoretical and numerical framework for the simultaneous identification of a spatially varying density coefficient and the initial displacement in a damped biharmonic wave equation. 
	By casting the forward system as an abstract Cauchy problem within a natural energy space, we utilized the Lumer--Phillips theorem to prove the well-posedness of the direct problem, demonstrating that the system's evolution constitutes a contraction semigroup. For the inverse problem, we derived a crucial energy observability inequality by employing the multiplier method under a star-shaped geometric control condition. 
	Building upon this observability estimate, we rigorously proved that both the heterogeneous density and the initial state can be simultaneously recovered from the boundary Cauchy data of the state's Laplacian with Lipschitz stability. A major mathematical contribution of this work is the explicit quantification of the damping mechanism's influence on the inverse problem. We analytically demonstrated that the Lipschitz stability constant scales precisely with the factor $(1 + \gamma)^{1/2}$. 
	Finally, our theoretical results were fully validated through numerical simulations in both one-dimensional and two-dimensional configurations. The numerical evidence confirms that while energy dissipation inherently reduces the signal-to-noise ratio at the boundary, the high-order structural properties of the biharmonic operator preserve sufficient information for robust parameter reconstruction, with the inversion error scaling predictably with respect to the damping coefficient.
	
\begin{thebibliography}{10}

\bibitem{Acosta2025Westervelt}
S.~Acosta and B.~Palacios.
\newblock Simultaneous determination of wave speed, diffusivity and
  nonlinearity in the westervelt equation using complex time-periodic
  solutions, sep 2025.
\newblock arXiv:2509.10718v1 [math.AP] 12 Sep 2025.

\bibitem{Arzhanykh1978}
I.~S. Ar\v{z}anyh.
\newblock Quasianalytic solutions of the biharmonic equation.
\newblock In {\em Direct and inverse problems for partial differential
  equations and their applications ({R}ussian)}, pages 55--61, 185. Izdat.
  ``Fan'' Uzbek. SSR, Tashkent, 1978.

\bibitem{Atahodzaev1974}
M.~A. Atahod\v{z}aev.
\newblock The exterior inverse biharmonic potential problem for a nearly
  spherical body.
\newblock In {\em Boundary value problems for differential equations, 4
  ({R}ussian)}, pages 77--93, 191. Izdat. ``Fan'' Uzbek. SSR, Tashkent, 1974.

\bibitem{Benrabah2019}
A.~Benrabah and N.~Boussetila.
\newblock Modified nonlocal boundary value problem method for an ill-posed
  problem for the biharmonic equation.
\newblock {\em Inverse Probl. Sci. Eng.}, 27(3):340--368, 2019.

\bibitem{Bhattacharyya2025}
S.~Bhattacharyya, K.~Krupchyk, S.~K. Sahoo, and G.~Uhlmann.
\newblock Inverse problems for third-order nonlinear perturbations of
  biharmonic operators.
\newblock {\em Comm. Partial Differential Equations}, 50(3):407--440, 2025.

\bibitem{MR3908}
D.~G. Bourgin.
\newblock The {D}irichlet problem for the damped wave equation.
\newblock {\em Duke Math. J.}, 7:97--120, 1940.

\bibitem{Bouslah2024}
Z.~Bouslah, A.~Hadj, and H.~Saker.
\newblock Shape reconstruction for an inverse biharmonic problem from partial
  {C}auchy data.
\newblock {\em Math. Methods Appl. Sci.}, 47(3):1613--1627, 2024.

\bibitem{MR4848562}
E.~E. Chitorkin and N.~P. Bondarenko.
\newblock Inverse {S}turm-{L}iouville problem with singular potential and
  spectral parameter in the boundary conditions.
\newblock {\em J. Differential Equations}, 421:495--523, 2025.

\bibitem{Danh2019}
H.~Q.~N. Danh, D.~O'Regan, V.~A. Vo, B.~T. Tran, and C.~H. Nguyen.
\newblock Regularization of an initial inverse problem for a biharmonic
  equation.
\newblock {\em Adv. Difference Equ.}, pages Paper No. 255, 20, 2019.

\bibitem{dong2024novel}
H.~Dong and P.~Li.
\newblock A novel boundary integral formulation for the biharmonic wave
  scattering problem.
\newblock {\em J. Sci. Comput.}, 98(2):Paper No. 42, 29, 2024.

\bibitem{dong2024uniqueness}
H.~Dong and P.~Li.
\newblock Uniqueness of an inverse cavity scattering problem for the
  time-harmonic biharmonic wave equation.
\newblock {\em Inverse Problems}, 40(6):Paper No. 065011, 19, 2024.

\bibitem{MR4959941}
A.~Feizmohammadi.
\newblock Reconstruction of 1{D} evolution equations and their initial data
  from one passive measurement.
\newblock {\em SIAM J. Math. Anal.}, 57(5):5089--5106, 2025.

\bibitem{MR21214}
F.~G. Friedlander.
\newblock Simple progressive solutions of the wave equation.
\newblock {\em Proc. Cambridge Philos. Soc.}, 43:360--373, 1947.

\bibitem{MR4598840}
Y.~Gao, H.~Liu, and Y.~Liu.
\newblock On an inverse problem for the plate equation with passive
  measurement.
\newblock {\em SIAM J. Appl. Math.}, 83(3):1196--1214, 2023.

\bibitem{Gazzola2010Polyharmonic}
F.~Gazzola, H.-C. Grunau, and G.~Sweers.
\newblock {\em Polyharmonic boundary value problems}, volume 1991 of {\em
  Lecture Notes in Mathematics}.
\newblock Springer-Verlag, Berlin, 2010.
\newblock Positivity preserving and nonlinear higher order elliptic equations
  in bounded domains.

\bibitem{harris2025direct}
I.~Harris, H.~Lee, and P.~Li.
\newblock Direct sampling for recovering a clamped cavity from the biharmonic
  far-field data.
\newblock {\em Inverse Problems}, 41(3):Paper No. 035013, 20, 2025.

\bibitem{li2026sampling}
I.~Harris, P.~Li, and G.~Ozochiawaeze.
\newblock Sampling methods for the inverse cavity scattering problem of
  biharmonic waves.
\newblock {\em Inverse Problems}, 42(1):Paper No. 015002, 23, 2026.

\bibitem{MR28496}
P.~Hartman and A.~Wintner.
\newblock A criterion for the non-degeneracy of the wave equation.
\newblock {\em Amer. J. Math.}, 71:206--213, 1949.

\bibitem{Kian2025WaveEquation}
Y.~Kian and H.~Liu.
\newblock Uniqueness and stability in determining the wave equation from a
  single passive boundary measurement, jul 2025.
\newblock arXiv:2507.10012v1 [math.AP] 14 Jul 2025.

\bibitem{Faouzi2025}
Y.~Kian and F.~Triki.
\newblock Recovery of an inclusion in photoacoustic imaging.
\newblock {\em Inverse Probl. Imaging}, 19(4):693--714, 2025.

\bibitem{li2026stability}
P.~Li, Z.~Li, and Y.~Liang.
\newblock Stability for inverse random source problems of the polyharmonic wave
  equation.
\newblock {\em Inverse Probl. Imaging}, 21:34--47, 2026.

\bibitem{li2025nonradiating}
P.~Li and J.~Wang.
\newblock Nonradiating sources of the biharmonic wave equation.
\newblock {\em Inverse Probl. Imaging}, 19(5):829--857, 2025.

\bibitem{Li2022a}
P.~Li and X.~Wang.
\newblock An inverse random source problem for the biharmonic wave equation.
\newblock {\em SIAM/ASA J. Uncertain. Quantif.}, 10(3):949--974, 2022.

\bibitem{Lions1988Controllability}
J.-L. Lions.
\newblock {\em Contr\^{o}labilit\'{e} exacte, perturbations et stabilisation de
  syst\`emes distribu\'{e}s. {T}ome 2}, volume~9 of {\em Recherches en
  Math\'{e}matiques Appliqu\'{e}es [Research in Applied Mathematics]}.
\newblock Masson, Paris, 1988.
\newblock Perturbations. [Perturbations].

\bibitem{Liu2025FourthOrderSchrodinger}
Y.~Liu and Y.~Gao.
\newblock Stability of inverse boundary value problem for the fourth-order
  schr{\"o}dinger equation.
\newblock {\em J. Math. Anal. Appl.} 2026, to appear.

\bibitem{MR1577103}
W.~H. McCrea and R.~A. Newing.
\newblock Boundary {C}onditions for the {W}ave {E}quation.
\newblock {\em Proc. London Math. Soc. (2)}, 37:520--534, 1934.

\bibitem{MR1577376}
H.~J. Priestley.
\newblock On {S}ome {S}olution of the {W}ave {E}quation.
\newblock {\em Proc. London Math. Soc. (2)}, 20(1):37--50, 1921.

\bibitem{Ramm1988}
A.~G. Ramm.
\newblock An inverse problem for biharmonic equation.
\newblock {\em Internat. J. Math. Math. Sci.}, 11(2):413--415, 1988.

\bibitem{wang2025adaptive}
Y.~Wang, P.~Li, L.~Xu, and T.~Yin.
\newblock An adaptive {D}irichlet-to-{N}eumann finite element method for the
  thermoelastic scattering problem.
\newblock {\em J. Comput. Phys.}, 534:Paper No. 114016, 24, 2025.

\bibitem{Zhao2025}
X.~Zhao and G.~Yuan.
\newblock Linearized inverse problem for biharmonic operators at high
  frequencies.
\newblock {\em Math. Methods Appl. Sci.}, 48(1):852--869, 2025.

\end{thebibliography}

	
	
	
	
	
	
	%\bibliographystyle{abbrv}
%	%	%%\bibliographystyle{elsarticle-num}
%	%	%%
%	%	%%%\bibliographystyle{UNRT}
%	%	%%%\bibliographystyle{IEEEtran}
%	%	%%\bibliographystyle{IEEEabrv,REFS}
%	%	%%%\bibliography{spbasic}
%	\bibliography{BGLM1}
\end{document} 
